\chapter{Estimating the theory from data}

\section{Estimate \(\G_i\) at the population-pair level}

A multivariate animal model or breeding design can estimate posterior draws
\begin{equation}
\G_i^{(s)},
\qquad s=1,\ldots,S.
\end{equation}
Population-pair-specific matrices are preferable because ecotype pooling can average over precisely the among-pair channel variation the theory seeks to explain.

\section{Estimate local curvature from phenotype--fitness data}

Near an optimum, fit a regularized quadratic selection model
\begin{equation}
\log W
=\alpha+\betav^\top\mathbf z
-\half\mathbf z^\top\Hh\mathbf z+\varepsilon.
\label{eq:empiricalQuadraticSelection}
\end{equation}
If traits are centred at the estimated optimum, the linear term should be small. If they are not centred, the optimum implied by the quadratic form is
\begin{equation}
\thetav=\Hh^{-1}\betav,
\end{equation}
provided \(\Hh\succ0\). This follows because
\begin{align}
\betav^\top\mathbf z-\half\mathbf z^\top\Hh\mathbf z
&=-\half(\mathbf z-\Hh^{-1}\betav)^\top\Hh
(\mathbf z-\Hh^{-1}\betav)
+\half\betav^\top\Hh^{-1}\betav.
\end{align}

\begin{cautionbox}
Quadratic regression is parameter hungry: \(n(n+1)/2\) curvature parameters are required. Shrinkage, structured low-rank models, preregistered trait modules, or Bayesian regularization are likely necessary. Positive-definiteness should be imposed or carefully assessed rather than assumed from noisy estimates.
\end{cautionbox}

\section{Construct posterior channel quantities}

For each posterior draw:
\begin{enumerate}[leftmargin=1.6em]
\item compute \(\G_i^{1/2}\);
\item compute \(\Sop_i=\G_i^{1/2}\Hh_i\G_i^{1/2}\);
\item choose or average over channel dimension \(k\);
\item map channel bases to phenotype space;
\item construct common-metric bases \(\Qq_i\);
\item compute principal angles, \(d_C^2\), \(I_V\), and spectral gaps;
\item compute exact expected log load and Gaussian arithmetic mean fitness for hybrids;
\item retain the full posterior joint distribution rather than plugging in posterior means.
\end{enumerate}

\section{Estimate hybrid moments}

For each cross and environment, estimate
\begin{equation}
\muv_H,
\qquad
\Pp_H.
\end{equation}
Then evaluate
\begin{equation}
L^{\log}_{H,e}
=\half(\muv_H-\thetav_e)^\top\Hh_e(\muv_H-\thetav_e)
+\half\tr(\Hh_e\Pp_H)
\end{equation}
and, where Gaussian approximation is defensible, \cref{eq:exactGaussianMeanFitness}.

\section{Competing predictive models}

A preregistered comparison should include:
\begin{enumerate}[leftmargin=1.6em]
\item divergence time or neutral genetic distance alone;
\item environmental distance alone;
\item mean mismatch alone;
\item \(\G\)-matrix distance alone;
\item \(\Hh\)-matrix distance alone;
\item channel angle or chordal distance alone;
\item the full curvature-weighted hybrid-load model;
\item the coupled ecological-plus-DMI model for later generations.
\end{enumerate}

The key evidentiary standard is out-of-sample prediction of hybrid fitness, not merely a significant association in the fitted data.
